\fontsize{9pt}{12pt}\selectfont

% Tiêu đề mục BÀI TẬP (EXERCISES)
{\fontsize{14pt}{16pt}\selectfont\textbf{\textsf{\color{stewartcyan}BÀI TẬP}}}\\[3pt]
{\color{stewartcyan}\rule{\linewidth}{1.5pt}}\\[9pt]

% Bố cục 2 cột chính tương xứng hoàn hảo với bản gốc
\noindent
\begin{minipage}[t]{0.485\textwidth}
  {\color{stewartcyan}\textbf{1--54}}\quad Tính $y'$.\\[6pt]
  % 1, 2
  \parbox[t]{0.48\linewidth}{\textbf{1.} $y = (x^2 + x^3)^4$}\hfill
  \parbox[t]{0.50\linewidth}{\textbf{2.} $y = \dfrac{1}{\sqrt{x}} - \dfrac{1}{\sqrt[5]{x^3}}}$\\[5.5pt]
  % 3, 4
  \parbox[t]{0.48\linewidth}{\textbf{3.} $y = \dfrac{x^2 - x + 2}{\sqrt{x}}}$}\hfill
  \parbox[t]{0.50\linewidth}{\textbf{4.} $y = \dfrac{\tan x}{1 + \cos x}$}\\[5.5pt]
  % 5, 6
  \parbox[t]{0.48\linewidth}{\textbf{5.} $y = x^2 \sin \pi x$}\hfill
  \parbox[t]{0.50\linewidth}{\textbf{6.} $y = x \cos^{-1} x$}\\[5.5pt]
  % 7, 8
  \parbox[t]{0.48\linewidth}{\textbf{7.} $y = \dfrac{t^4 - 1}{t^4 + 1}$}\hfill
  \parbox[t]{0.50\linewidth}{\textbf{8.} $x e^y = y \sin x$}\\[5.5pt]
  % 9, 10
  \parbox[t]{0.48\linewidth}{\textbf{9.} $y = \ln(x \ln x)$}\hfill
  \parbox[t]{0.50\linewidth}{\textbf{10.} $y = e^{mx} \cos nx$}\\[5.5pt]
  % 11, 12
  \parbox[t]{0.48\linewidth}{\textbf{11.} $y = \sqrt{x} \cos \sqrt{x}$}\hfill
  \parbox[t]{0.50\linewidth}{\textbf{12.} $y = (\arcsin 2x)^2$}\\[5.5pt]
  % 13, 14
  \parbox[t]{0.48\linewidth}{\textbf{13.} $y = \dfrac{e^{1/x}}{x^2}$}\hfill
  \parbox[t]{0.50\linewidth}{\textbf{14.} $y = \ln \sec x$}\\[5.5pt]
  % 15, 16
  \parbox[t]{0.48\linewidth}{\textbf{15.} $y + x \cos y = x^2 y$}\hfill
  \parbox[t]{0.50\linewidth}{\textbf{16.} $y = \left(\dfrac{u - 1}{u^2 + u + 1}\right)^4$}\\[5.5pt]
  % 17, 18
  \parbox[t]{0.48\linewidth}{\textbf{17.} $y = \sqrt{\arctan x}$}\hfill
  \parbox[t]{0.50\linewidth}{\textbf{18.} $y = \cot(\csc x)$}\\[5.5pt]
  % 19, 20
  \parbox[t]{0.48\linewidth}{\textbf{19.} $y = \tan\left(\dfrac{t}{1 + t^2}\right)$}\hfill
  \parbox[t]{0.50\linewidth}{\textbf{20.} $y = e^{x \sec x}$}\\[5.5pt]
  % 21, 22
  \parbox[t]{0.48\linewidth}{\textbf{21.} $y = 3^{x \ln x}$}\hfill
  \parbox[t]{0.50\linewidth}{\textbf{22.} $y = \sec(1 + x^2)$}\\[5.5pt]
  % 23, 24
  \parbox[t]{0.48\linewidth}{\textbf{23.} $y = (1 - x^{-1})^{-1}$}\hfill
  \parbox[t]{0.50\linewidth}{\textbf{24.} $y = 1 / \sqrt[3]{x + \sqrt{x}}$}\\[5.5pt]
  % 25, 26
  \parbox[t]{0.48\linewidth}{\textbf{25.} $\sin(xy) = x^2 - y$}\hfill
  \parbox[t]{0.50\linewidth}{\textbf{26.} $y = \sqrt{\sin \sqrt{x}}$}\\[5.5pt]
  % 27, 28
  \parbox[t]{0.48\linewidth}{\textbf{27.} $y = \log_5(1 + 2x)$}\hfill
  \parbox[t]{0.50\linewidth}{\textbf{28.} $y = (\cos x)^x$}\\[5.5pt]
  % 29, 30
  \parbox[t]{0.48\linewidth}{\textbf{29.} $y = \ln \sin x - \frac{1}{2} \sin^2 x$}\hfill
  \parbox[t]{0.50\linewidth}{\textbf{30.} $y = \dfrac{(x^2 + 1)^4}{(2x + 1)^3(3x - 1)^5}$}\\[5.5pt]
  % 31, 32
  \parbox[t]{0.48\linewidth}{\textbf{31.} $y = x \tan^{-1}(4x)$}\hfill
  \parbox[t]{0.50\linewidth}{\textbf{32.} $y = e^{\cos x} + \cos(e^x)$}\\[5.5pt]
  % 33, 34
  \parbox[t]{0.48\linewidth}{\textbf{33.} $y = \ln|\sec 5x + \tan 5x|$}\hfill
  \parbox[t]{0.50\linewidth}{\textbf{34.} $y = 10^{\tan \pi\theta}$}\\[5.5pt]
  % 35, 36
  \parbox[t]{0.48\linewidth}{\textbf{35.} $y = \cot(3x^2 + 5)$}\hfill
  \parbox[t]{0.50\linewidth}{\textbf{36.} $y = \sqrt{t \ln(t^4)}$}\\[5.5pt]
  % 37, 38
  \parbox[t]{0.48\linewidth}{\textbf{37.} $y = \sin(\tan \sqrt{1 + x^3})$}\hfill
  \parbox[t]{0.50\linewidth}{\textbf{38.} $y = x \sec^{-1} x$}\\[5.5pt]
  % 39
  \textbf{39.} $y = 5^{\arctan(1/x)}$\\[5.5pt]
  % 40
  \textbf{40.} $y = \sin^{-1}(\cos\theta), \quad 0 < \theta < \pi$\\[5.5pt]
  % 41
  \textbf{41.} $y = x \tan^{-1} x - \frac{1}{2}\ln(1 + x^2)$\\[5.5pt]
  % 42, 43
  \parbox[t]{0.48\linewidth}{\textbf{42.} $y = \ln(\arcsin x^2)$}\hfill
  \parbox[t]{0.50\linewidth}{\textbf{43.} $y = \tan^2(\sin\theta)$}\\[5.5pt]
  % 44, 45
  \parbox[t]{0.42\linewidth}{\textbf{44.} $y + \ln y = xy^2$}\hfill
  \parbox[t]{0.56\linewidth}{\textbf{45.} $y = \dfrac{\sqrt{x + 1}\,(2 - x)^5}{(x + 3)^7}$}
\end{minipage}%
\hfill
\begin{minipage}[t]{0.485\textwidth}
  % 46, 47
  \parbox[t]{0.48\linewidth}{\textbf{46.} $y = \dfrac{(x + \lambda)^4}{x^4 + \lambda^4}$}\hfill
  \parbox[t]{0.50\linewidth}{\textbf{47.} $y = x \sinh(x^2)$}\\[5.5pt]
  % 48, 49
  \parbox[t]{0.48\linewidth}{\textbf{48.} $y = \dfrac{\sin mx}{x}$}\hfill
  \parbox[t]{0.50\linewidth}{\textbf{49.} $y = \ln(\cosh 3x)$}\\[5.5pt]
  % 50, 51
  \parbox[t]{0.48\linewidth}{\textbf{50.} $y = \ln\sqrt{\dfrac{x^2 - 4}{2x + 5}}$}\hfill
  \parbox[t]{0.50\linewidth}{\textbf{51.} $y = \cosh^{-1}(\sinh x)$}\\[5.5pt]
  % 52, 53
  \parbox[t]{0.48\linewidth}{\textbf{52.} $y = x \tanh^{-1}\sqrt{x}$}\hfill
  \parbox[t]{0.50\linewidth}{\textbf{53.} $y = \cos(e^{\sqrt{\tan 3x}})}$}\\[5.5pt]
  % 54
  \textbf{54.} $y = \sin^2(\cos\sqrt{\sin\pi x})$\\[11pt]

  % 55-60
  \textbf{55.} Cho $f(t) = \sqrt{4t + 1}$, hãy tìm $f''(2)$.\\[4pt]
  \textbf{56.} Cho $g(\theta) = \theta \sin\theta$, hãy tìm $g''(\pi/6)$.\\[4pt]
  \textbf{57.} Tìm $y''$ nếu $x^6 + y^6 = 1$.\\[4pt]
  \textbf{58.} Tìm $f^{(n)}(x)$ nếu $f(x) = 1/(2 - x)$.\\[4pt]
  \textbf{59.} Dùng phương pháp quy nạp toán học để chứng minh rằng nếu $f(x) = x e^x$, thì $f^{(n)}(x) = (x + n)e^x$. (\textit{Ghi chú:} Xem Các nguyên lý giải toán sau Chương 1.)\\[4pt]
  \textbf{60.} Tính giới hạn $\displaystyle\lim_{t \to 0}\frac{t^3}{\tan^3(2t)}$.\\[11pt]

  % 61-63
  {\color{stewartcyan}\textbf{61--63}}\quad Tìm phương trình tiếp tuyến của đường cong tại điểm đã cho.\\[5pt]
  \parbox[t]{0.50\linewidth}{\textbf{61.} $y = 4 \sin^2 x, \quad (\pi/6, 1)$}\hfill
  \parbox[t]{0.48\linewidth}{\textbf{62.} $y = \dfrac{x^2 - 1}{x^2 + 1}, \quad (0, -1)$}\\[5pt]
  \textbf{63.} $y = \sqrt{1 + 4 \sin x}, \quad (0, 1)$\\[11pt]

  % 64-65
  {\color{stewartcyan}\textbf{64--65}}\quad Tìm phương trình tiếp tuyến và pháp tuyến của đường cong tại điểm đã cho.\\[5pt]
  \textbf{64.} $x^2 + 4xy + y^2 = 13, \quad (2, 1)$\\[4.5pt]
  \textbf{65.} $y = (2 + x)e^{-x}, \quad (0, 2)$\\[11pt]

  % 66
  \noindent\llap{\casicon\hspace{4pt}}\textbf{66.} Cho $f(x) = x e^{\sin x}$, hãy tìm $f'(x)$. Vẽ đồ thị của $f$ và $f'$ trên cùng một màn hình và nhận xét.\\[7pt]

  % 67
  \textbf{67.} (a) Cho $f(x) = x\sqrt{5 - x}$, hãy tìm $f'(x)$.\\[3pt]
  \hspace*{1.4em}(b) Tìm phương trình các tiếp tuyến của đường cong $y = x\sqrt{5 - x}$ tại các điểm $(1, 2)$ và $(4, 4)$.\\[3pt]
  \noindent\llap{\casicon\hspace{4pt}}\hspace*{1.4em}(c) Minh họa câu (b) bằng cách vẽ đồ thị của đường cong và các tiếp tuyến trên cùng một màn hình.\\[3pt]
  \noindent\llap{\casicon\hspace{4pt}}\hspace*{1.4em}(d) Kiểm tra xem câu trả lời ở câu (a) có hợp lý hay không bằng cách so sánh đồ thị của $f$ và $f'$.
\end{minipage}
